\chapter{Conclusion}
\section{Conclusion}
This project set out to treat the KV cache of a transformer language model, its runtime memory, as an object that can be measured, attributed and edited, and to use social bias in two-agent conversations as the test case. The motivation comes from SDG~10: as language models are combined into systems of cooperating agents, bias that enters one agent could spread to others through shared memory without appearing in any transcript, and the people affected by such systems are mostly those already disadvantaged.

The implemented contribution is an open Python framework, together with its evaluation and a public demonstration. The framework stamps every cache entry of a frozen model with its layer, head, position, turn, agent and channel; runs two-agent conversations in which information crosses through transplanted cache entries, through text, or not at all; detects persistent memory attractors while holding out attention sinks; localises bias reservoirs with cross-validated probes and false-discovery control; separates memory-mediated from text-mediated bias transfer; tests for phase transitions under strict criteria; and evaluates cache-only interventions by seed-matched replay against random-head and sink-only controls and a pre-registered coherence budget. It is covered by 36 automated tests, reproduces its results exactly, and ran 270 conversations on a laptop, including a real, publicly released model on a 4~GB GPU.

The evidence supports the following conclusions, which are consistent with Chapter~6. Runtime memory determines the answer of both models: six one-sided sentences moved Qwen2.5-0.5B from almost never to almost always choosing the stereotype trait. Bias reached the second agent in both models. In the reference model, shared cache entries carried more bias than text, with memory contagion of $+0.186$ and a 95\% interval of $+0.034$ to $+0.340$, positive at all nine pressure levels; the rank-based test gave $p = 0.085$, so this result is consistent but not yet conclusive. In Qwen2.5 the text carried the bias and the cache added nothing measurable, because the sender repeated the statements in its replies. Transplanted entries became persistent attractors in the reference model, and middle-layer heads predicted bias with AUC up to 0.92, while Qwen2.5 showed no persistent structure beyond the attention sink. No phase transition was found in either model. The strongest intervention, ablation of the entries read by the reservoir heads, reduced bias by 20.3\% and beat both controls, but it cost 6.7 points of factual accuracy and was correctly rejected.

Measured against the objectives, O1 (instrumentation) and O4 (contagion and phase analysis) were achieved; O2 (attractors) was achieved in the reference model, with a negative finding in Qwen2.5; O3 (reservoirs) was partly achieved, with strong predictive heads but small single-head causal effects; and O5 (causal control) was achieved as a method, with every component validated, but no intervention met the coherence budget. Taken together, the project shows that the memory channel between agents can carry bias, that whether it does depends on the model and on how much the exchanged text already carries, and that cache edits can reduce bias but, in the setting studied, not without damaging other knowledge. These are the questions the framework was built to answer, and it answers them reproducibly on commodity hardware.

\section{Future work}
The work planned for the remaining part of the project (Semester~8) follows the timeline of the design document and addresses the limitations of Section~6.6.

The first and most important task is to \textbf{run the study on 7--8 billion parameter models}, namely Llama-3-8B-Instruct, Mistral-7B-Instruct and Qwen2.5-7B-Instruct, for which configuration files are already provided. This depends on access to a GPU with at least 24~GB of memory, through the institute cluster or a rented cloud instance, and on acceptance of the Llama~3 licence. Validation consists of repeating the complete pipeline unchanged and comparing contagion, attractor and reservoir results across the three models, which also makes the cross-model replication criterion of the phase analysis testable.

The second task is to \textbf{strengthen the statistical evidence for memory contagion} by running more seeds and longer conversations. The reference result has an interval excluding zero but a Mann--Whitney $p$ of 0.085 with five seeds per cell; the number of seeds needed for a conclusive test will be fixed in advance from the observed effect size, and longer conversations will also give per-turn bias series long enough for a meaningful transfer-entropy estimate.

The third task is to \textbf{broaden and harden the evaluation}: add probes built from the BBQ, StereoSet, CrowS-Pairs and TruthfulQA datasets, using real group terms alongside the fictional controls so that learned and contextual bias can be separated; measure the coherence budget on MMLU and GSM8K subsets as specified in the design document; and measure the instrumentation overhead against uninstrumented generation, with a target of less than twice the uninstrumented time.

The fourth task is to \textbf{refine the interventions} so that they remove bias without removing factual knowledge. Since the rejected ablation beat both controls but cost accuracy, the next step is to select target entries that predict bias but not factual answers, for example by excluding entries that the factual probes also depend on, and to sweep the perturbation strength and gate factor under the same acceptance rule. Every attempt will continue to be reported, including rejected ones.

The final tasks are the planned deliverables: public release of the trace dataset with a data card and of the toolkit under an open licence, submission of a research paper describing the framework and the contagion results, the final dissertation, and the final demonstration and viva in April~2027. The demonstration website will be extended with the results for the larger models as they become available, keeping the website, the presentation and this report consistent.
